% Generated draft: requires mathematical and source review.
\documentclass[11pt]{article}
\usepackage[margin=25mm]{geometry}
\usepackage{fontspec}
\setmainfont{Latin Modern Roman}
\usepackage{amsmath,amssymb,mathtools,unicode-math}
\setmathfont{Latin Modern Math}
\usepackage{enumitem,needspace,xurl,hyperref}
\setlength{\parindent}{0pt}
\setlength{\parskip}{5pt}
\setlength{\emergencystretch}{3em}
\newcommand{\N}{\mathbb{N}}
\newcommand{\Z}{\mathbb{Z}}
\newcommand{\Q}{\mathbb{Q}}
\newcommand{\R}{\mathbb{R}}
\newcommand{\C}{\mathbb{C}}
\newcommand{\F}{\mathbb{F}}
\newcommand{\A}{\mathbb{A}}
\newcommand{\PP}{\mathbb{P}}
\newcommand{\E}{\mathbb{E}}
\newcommand{\eps}{\varepsilon}
\newcommand{\dd}{\,\mathrm{d}}
\begin{document}
\section*{UnsolvedMath ID: 30005697}
\textbf{Source problem code:} catalogue entry
\subsection{Definitions and mathematical statement}
Write \(\mathcal A_G\) for the category denoted by \(A_G\) in the supplied record. Its definition, the allowed groups \(G\), any coefficient data, and the ambient hypotheses of the source's Theorem C are not reproduced. Consequently, the following is a conditional formulation within that source setting, rather than a statement about arbitrary groups or arbitrary abelian categories.

Assuming the usual abelian-category setting with enough projectives and injectives, let \(\operatorname{Proj}(\mathcal A_G)\) and \(\operatorname{Inj}(\mathcal A_G)\) denote its projective and injective objects. With the standard meanings of the invariant names, set
\[
\operatorname{silp}(\mathcal A_G)=\sup\{\operatorname{id}_{\mathcal A_G}P:P\in\operatorname{Proj}(\mathcal A_G)\},\qquad
\operatorname{spli}(\mathcal A_G)=\sup\{\operatorname{pd}_{\mathcal A_G}I:I\in\operatorname{Inj}(\mathcal A_G)\}.
\]
Here \(\operatorname{id}\) and \(\operatorname{pd}\) are injective and projective dimension, allowing infinite values. These conventions and the required existence of resolutions must be checked against the missing source context.

Let \(D^b(\mathcal A_G)\) be the bounded derived category and let \(K^b(\operatorname{Proj}(\mathcal A_G))\) and \(K^b(\operatorname{Inj}(\mathcal A_G))\) be the homotopy categories of bounded complexes of projectives and injectives, viewed inside the derived category. The two displayed quotient categories in the source fragment are the Verdier quotients
\[
D^b(\mathcal A_G)/K^b(\operatorname{Proj}(\mathcal A_G)),\qquad
D^b(\mathcal A_G)/K^b(\operatorname{Inj}(\mathcal A_G)).
\]
A totally acyclic complex of projectives is an exact complex \(P^\bullet\) of projective objects for which \(\operatorname{Hom}_{\mathcal A_G}(P^\bullet,Q)\) is exact for every projective \(Q\). Dually, a totally acyclic complex of injectives is an exact complex \(I^\bullet\) of injective objects for which \(\operatorname{Hom}_{\mathcal A_G}(J,I^\bullet)\) is exact for every injective \(J\). Denote their homotopy categories by \(K_{\mathrm{tac}}(\operatorname{Proj}(\mathcal A_G))\) and \(K_{\mathrm{tac}}(\operatorname{Inj}(\mathcal A_G))\).

Gorenstein-projective objects are cycles of totally acyclic complexes of projectives; Gorenstein-injective objects are defined dually. Their stable categories, denoted \(\underline{\operatorname{GProj}}(\mathcal A_G)\) and \(\underline{\operatorname{GInj}}(\mathcal A_G)\), retain these objects and quotient morphisms by those factoring through projectives and injectives, respectively. The supplied fragment explicitly describes the Gorenstein-projective category as Frobenius with projective-injective objects precisely the \(\mathcal A_G\)-projectives.

The fragment also describes a complete-cohomology category having all objects of \(\mathcal A_G\) as objects and zeroth complete cohomology, computed using injective resolutions, as its morphism groups. Its complete-cohomology construction and the first item of the source's category list are missing.

\textbf{Question, conditional on the unreproduced source hypotheses.} For every instance of \(\mathcal A_G\) allowed by Theorem C, does that theorem's equivalence conclusion still hold under either one of the separate assumptions
\[
\operatorname{silp}(\mathcal A_G)<\infty
\quad\text{or}\quad
\operatorname{spli}(\mathcal A_G)<\infty,
\]
without assuming finiteness of the other invariant? The intended conclusion concerns equivalences among the complete-cohomology categories, the two Verdier quotients above, the two homotopy categories of totally acyclic complexes, and the stable Gorenstein-projective and Gorenstein-injective categories. The supplied curation reports that Theorem C establishes its conclusion when both invariants are finite. The exact full category list, comparison functors, and structure preserved by the equivalences cannot be recovered from the supplied fragment; this is therefore a statement-only draft with unresolved source qualifications.
\subsection{Short English statement}
In the setting of the source's Theorem C, ask whether finiteness of either \(\operatorname{silp}(\mathcal A_G)\) alone or \(\operatorname{spli}(\mathcal A_G)\) alone suffices for its equivalences between complete-cohomology, derived quotient, totally acyclic complex, and stable Gorenstein categories. The supplied record omits the definition of \(\mathcal A_G\), the full hypotheses, and part of the conclusion, so a fully specified theorem statement requires source review.
\subsection{Sources}
\begin{itemize}
\item Supplied local UnsolvedMath registry record, permanent ID \(30005697\), title \emph{Gorenstein Equivalences under One-Sided Finiteness}; source problem code \texttt{catalogue entry}, in the Oberwolfach Reports Open Problems collection. No individual proposer is supplied.
\item Supplied citation: \emph{Mini-Workshop: Homological aspects of TDLC-groups (2024)}, Oberwolfach Reports, pp.~3103--3138, DOI \texttt{10.4171/owr/2023/55}. \url{https://doi.org/10.4171/owr/2023/55}. The record separately gives proposed year 2023; the relationship between these year labels has not been verified.
\item Supplied literature source, labelled ``Oberwolfach report'': \url{https://ems.press/content/serial-article-files/48166}.
\item Literature assessment dated exactly \texttt{2026-08-22}: the supplied curation labels the problem open and reports that Theorem C is proved assuming both invariants finite, while its ending question asks about one-sided finiteness. This is supplied dated curation, not independently verified literature assessment.
\item The registry labels the statement \texttt{corrected\_verified} and reports that it was checked against the official report. Those are supplied verification claims; no independent source consultation or mathematical verification is asserted here.
\end{itemize}
\end{document}
